% Einstiegsblatt, erzeugt aus einstieg/erzeugen.py – setzen mit ../bauen.sh
\documentclass[a4paper, 9.5pt, parskip=half]{scrartcl}
\usepackage[dutch]{babel}
\usepackage[margin=16mm, top=14mm, bottom=14mm]{geometry}
\usepackage{wtwinci}
\usepackage{csquotes}
\usepackage{multicol}
\graphicspath{{/home/thobug/AI-Projects/app4webtrees/docs/screenshots/}{/home/thobug/AI-Projects/app4webtrees/docs/}}
\pagestyle{empty}
\setlength{\columnsep}{8mm}
\begin{document}
\noindent\begin{minipage}[c]{0.13\linewidth}\includegraphics[width=\linewidth]{icon/icon-512.png}\end{minipage}\hfill
\begin{minipage}[c]{0.84\linewidth}\sffamily{\fontsize{26}{30}\selectfont\bfseries\color{wtPetrol}wtWin in vijf minuten\par}\vspace{3pt}{\small\color{wtGrau}Uw webtrees-stamboom als desktopprogramma · wtWin (Windows) · wtTux (Linux) · wtMac (macOS, om te testen)\par}\end{minipage}
\vspace{4mm}\noindent{\color{wtPetrol}\rule{\linewidth}{0.8pt}}\vspace{2mm}
\begin{multicols}{2}
\section*{1 · Installeren}
\textbf{Windows:} download \texttt{wtWin-….exe} van \href{https://wtwin.bgg-home.de}{wtwin.bgg-home.de} en start het. Het installatieprogramma vraagt eerst de taal, heeft geen beheerdersrechten nodig en vervangt een oudere versie vanzelf. Bij de eerste start: „Meer informatie“ › „Toch uitvoeren“ (niet ondertekend).

\textbf{Linux:} \texttt{sudo apt install ./wttux\_…\_amd64.deb}

\textbf{macOS:} \texttt{…-arm64.dmg} (Apple-chip) of \texttt{…-x64.dmg} (Intel); bij de eerste start „Toch openen“ onder Systeeminstellingen › Privacy.
\section*{2 · Verbinden of alleen werken}
\textbf{Met uw webtrees:} meld u in de browser aan bij webtrees, menu \menue{App} › \menue{Verbinden met wtWin}. Het programma neemt adres en account over, vraagt één keer en opent de stamboom. Vereist: webtrees 2.2 met de module api4webtrees (zit in nas4webtrees).

\textbf{Handmatig:} voer het adres van uw webtrees-site in, \menue{Verbinden}, meld u aan met gebruikersnaam en wachtwoord. Dezelfde rechten als op de website gelden.

\textbf{Zonder server:} rechts op het startscherm \menue{Stamboom aanmaken} of \menue{Uit GEDCOM-bestand overnemen …}. Daarna werkt u zonder server, wachtwoord en internet; webtrees draait onzichtbaar op de achtergrond.
\section*{3 · Drie indelingen}
\textbf{Navigator} (\taste{Ctrl+1}): links de kinderen, in het midden de persoon met huwelijken, rechts vier generaties voorouders. Eén klik zet iemand anders in het midden.

\textbf{Boom in het midden} (\taste{Ctrl+2}): personenlijst, boom en persoonspaneel naast elkaar; volledig scherm om in de boom te werken.

\textbf{Gezinsweergave} (\taste{Ctrl+3}): voorouders van beide partners boven, het paar in het midden, kinderen met partners en kleinkinderen onder.
\section*{4 · Maken en controleren}
\textbf{Grafieken} (15 soorten, ook plotter en grootformaat), \textbf{lijsten} en \textbf{boeken} (voorouders, nakomelingen, familieboek, ook als website) in het menu \menue{Maken}, als PDF of afdruk.

\textbf{Extra:} plausibiliteitscontrole met 61 regels, personen samenvoegen, plaatsen opschonen, datumcalculator, kerkelijk jaar.

\textbf{Bronnen, plaatsen, archief, taken:} eigen vensters onder \menue{Beeld} (\taste{Ctrl+5} tot \taste{Ctrl+9}).
\section*{5 · Hulp}
\taste{F1} opent de hulp bij het actieve venster. \taste{Ctrl+F} zoekt een persoon, \taste{F5} laadt opnieuw, \taste{Ctrl+C} kopieert de persoonstekst.

Handleiding als PDF, vragen en fouten: \href{https://wtwin.bgg-home.de}{wtwin.bgg-home.de} · \href{https://github.com/thobgg/app4webtrees/issues}{github.com/thobgg/app4webtrees/issues}
\vspace{3mm}
\noindent\bild{desktop-navigator.png}
\end{multicols}
\vspace{2mm}
\noindent{\color{wtLinie}\rule{\linewidth}{0.5pt}}\\[2pt]
{\sffamily\footnotesize\color{wtGrau}wtWin 1.52 · GPL-3.0 · wtwin.bgg-home.de}
\end{document}
